% ============================================================
%  Izveštaj o projektu, Neuronske mreže (13E054NM)
%  Deo 1: Prepoznavanje emocija iz govora pomoću CNN-a
%  Deo 2: Fuzzy logika (biće dodat, videti marker na dnu)
%
%  Kompajlirati sa pdflatex (npr. na Overleaf-u).
%  Slike se učitavaju iz foldera results/, okačiti ceo folder.
% ============================================================
\documentclass[12pt,a4paper]{article}

\usepackage[utf8]{inputenc}
\usepackage[T1]{fontenc}
\usepackage[serbian]{babel}
\usepackage[margin=2.5cm]{geometry}
\usepackage{graphicx}
\graphicspath{{results/}}
\usepackage{booktabs}
\usepackage{amsmath}
\usepackage{float}
\usepackage[hidelinks]{hyperref}

\title{Prepoznavanje emocija iz govora\\
       \large Projekat iz predmeta Neuronske mreže (13E054NM)\\
       Elektrotehnički fakultet, Univerzitet u Beogradu}
% TODO: upisati ime(na) i broj(eve) indeksa
\author{Ime Prezime, broj indeksa}
\date{\today}

\begin{document}

\maketitle
\tableofcontents
\newpage

% ============================================================
\part{Prepoznavanje emocija iz govora pomoću CNN-a}
% ============================================================

\section{Opis problema}

Cilj projekta je automatsko prepoznavanje emocije iz snimka govora. Svaki
audio fajl pretvara se u mel-spektrogram, 2D sliku koja prikazuje
raspoređenost energije zvuka kroz vreme i frekvenciju, a zatim se ta slika
klasifikuje konvolucionom neuronskom mrežom u jednu od 8 klasa emocija.

Prepoznavanje emocija iz govora je inherentno težak problem jer:
\begin{itemize}
  \item ista emocija može zvučati različito kod različitih govornika;
  \item neke emocije (tuga, strah, neutralno) imaju slične akustičke
        karakteristike;
  \item skup podataka je mali (1.440 snimaka za 8 klasa), što pogoduje
        preprilagođavanju;
  \item ni ljudi nisu pouzdani u ovom zadatku: autori RAVDESS baze mere
        ljudsku tačnost od $\approx 62\%$ na audio snimcima govora, što je
        realan orijentir za plafon zadatka.
\end{itemize}

\section{Skup podataka: RAVDESS}

\begin{table}[H]
\centering
\begin{tabular}{ll}
\toprule
Karakteristika & Vrednost \\
\midrule
Ukupno snimaka           & 1.440 \\
Broj glumaca             & 24 (12 muških, 12 ženskih) \\
Broj klasa emocija       & 8 \\
Frekvencija uzorkovanja  & 48\,kHz (resample na 22.050\,Hz) \\
Format                   & WAV (mono) \\
\bottomrule
\end{tabular}
\caption{Osnovne karakteristike RAVDESS skupa podataka.}
\end{table}

\noindent\textbf{Klase emocija:} \texttt{neutral}, \texttt{calm},
\texttt{happy}, \texttt{sad}, \texttt{angry}, \texttt{fearful},
\texttt{disgust}, \texttt{surprised}.

\subsection{Neravnoteža klasa}

Klasa \texttt{neutral} ima samo 96 snimaka, dok svaka od preostalih 7 klasa
ima 192 snimka, odnos 2:1 (slika~\ref{fig:dist}). Problem je rešen
tačno balansiranom augmentacijom: neutral dobija duplo više augmentovanih
kopija, pa svih 8 klasa ulazi u trening sa identičnim brojem uzoraka (432).

\begin{figure}[H]
  \centering
  \includegraphics[width=0.85\textwidth]{class_distribution.png}
  \caption{Raspodela uzoraka po klasama emocija.}
  \label{fig:dist}
\end{figure}

\begin{figure}[H]
  \centering
  \includegraphics[width=\textwidth]{sample_spectrograms.png}
  \caption{Po jedan mel-spektrogram za svaku klasu emocija.}
  \label{fig:samples}
\end{figure}

\section{Preprocesiranje audio signala}

\subsection{Pipeline po fajlu}

\begin{enumerate}
  \item učitavanje celog snimka (librosa, $sr = 22.050$\,Hz);
  \item trim tišine ($\mathrm{top\_db}=30$);
  \item crop/padding na tačno 3 sekunde;
  \item [opciono] augmentacija talasnog oblika;
  \item mel-spektrogram (128 mel binova), zatim konverzija u dB skalu;
  \item delta i delta-delta kanali (brzina i ubrzanje spektra);
  \item resize svakog kanala na $128\times128$ (bilinearna interpolacija);
  \item Z-score normalizacija po kanalu;
  \item [opciono] SpecAugment (maske preko sva 3 kanala);
  \item izlaz: tenzor $(128, 128, 3)$.
\end{enumerate}

\subsection{Trim tišine}

RAVDESS snimci počinju i završavaju se tišinom. Merenje na tipičnom fajlu:
od 3,30\,s snimka, posle trima ostane 1,32\,s govora. Bez trima bi ulaz u
mrežu bio 40--60\% tišina, a kraj rečenice bi često bio odsečen jer se seku
prve 3 sekunde snimka. Funkcija \texttt{librosa.effects.trim(top\_db=30)}
uklanja delove tiše od 30\,dB ispod vrha snimka; prag je relativan, pa se
prilagođava i tihim (\texttt{sad}, \texttt{calm}) snimcima.

\subsection{Delta kanali}

Statični mel-spektrogram opisuje šta je prisutno u spektru, ali ne i kako
se spektar menja, a emocija je upravo u dinamici govora (tempo, tremor,
nagle promene). Zato se uz mel-spektrogram računaju delta (prvi izvod po
vremenu) i delta-delta (drugi izvod), pa se slažu kao tri kanala slike.
Dobitak je dvostruk:
\begin{enumerate}
  \item CNN od nule dobija eksplicitnu informaciju o dinamici govora;
  \item transfer model (EfficientNetB0 očekuje RGB ulaz) dobija tri
        informativna kanala umesto tri kopije istog kanala.
\end{enumerate}

\subsection{Z-score normalizacija po uzorku i kanalu}

Različiti snimci imaju različite nivoe glasnoće. Z-score normalizacijom
svaki kanal svakog spektrograma centrira se oko nule i skalira na jediničnu
devijaciju, nezavisno od glasnoće snimka, pa mreža ne može koristiti
apsolutnu glasnoću kao prečicu za klasifikaciju.

\subsection{Podela podataka}

\begin{table}[H]
\centering
\begin{tabular}{lrr}
\toprule
Skup & Udeo & Broj uzoraka (pre augmentacije) \\
\midrule
Trening    & 75\% & 1.080 \\
Validacija & 15\% & 216 \\
Test       & 10\% & 144 \\
\bottomrule
\end{tabular}
\caption{Stratifikovana nasumična podela (fiksni seed, deterministična).}
\end{table}

Test skup od 144 uzorka znači da jedna predikcija menja tačnost za
0,7\,p.p., pa su razlike manje od $\sim$3\,p.p.\ u zoni šuma.

\emph{Napomena:} tokom razvoja testirana je i speaker-independent podela
(glumci 21--24 izdvojeni u validaciju i test), ali daje $\sim$38\% tačnosti
jer se glumci akustički previše razlikuju da bi 2 test glumca bila
reprezentativna. Zadržana je nasumična stratifikovana (speaker-dependent)
podela, što treba imati u vidu pri poređenju sa literaturom.

\section{Augmentacija podataka}

\subsection{Augmentacija talasnog oblika}

Augmentacija se primenjuje na sirovi talasni oblik, ne na piksele
spektrograma: augmentacija piksela stvara fizički nerealne artefakte, dok
augmentacija talasnog oblika uvek daje validan audio signal.

\begin{table}[H]
\centering
\begin{tabular}{lll}
\toprule
Tehnika & Parametri & Efekat na spektrogram \\
\midrule
Time stretching & rate $\in [0.85, 1.15]$, $p{=}0.5$ & horizontalno širenje/skupljanje \\
Pitch shifting  & $\pm 2$ polutona, $p{=}0.5$        & vertikalni pomak sadržaja \\
Gaussov šum     & $\sigma{=}0.005$, $p{=}0.5$        & simulacija šuma mikrofona \\
\bottomrule
\end{tabular}
\caption{Augmentacije talasnog oblika.}
\end{table}

Svaka augmentovana kopija ima deterministički seed, pa je ceo trening skup
reprodukovljiv iako se preprocesiranje izvršava paralelno na svim jezgrima.

\subsection{SpecAugment}

SpecAugment maskira nasumične frekvencijske pojaseve i vremenske segmente
(2 maske $\times$ maksimalno 20 binova/okvira). Primenjuje se posle Z-score
normalizacije (maskirana vrednost 0 je približno sredina distribucije) i
preko sva tri kanala istovremeno, da mel i delta kanali ostanu poravnati.

\subsection{Tačno balansirana augmentacija}

\begin{table}[H]
\centering
\begin{tabular}{lrrr}
\toprule
Klasa & Original & Kopija (1 plain + N aug) & Ukupno \\
\midrule
neutral          & 72  & $6\times$ (1+5) & 432 \\
svih ostalih 7   & 144 & $3\times$ (1+2) & 432 \\
\bottomrule
\end{tabular}
\caption{Trening skup: 3.456 uzoraka, tačno 432 po klasi. Validacija i
test se nikad ne augmentuju.}
\end{table}

Pošto augmentacija već izjednačava brojeve uzoraka po klasama, class
weights nisu potrebni.

\section{Arhitektura modela: CNN od nule}

\subsection{Struktura}

\begin{center}
\begin{tabular}{ll}
\toprule
Sloj & Izlaz \\
\midrule
Ulaz (mel + $\Delta$ + $\Delta\Delta$) & $(128, 128, 3)$ \\
CNNBlock(32)                    & $(64, 64, 32)$ \\
CNNBlock(64)                    & $(32, 32, 64)$ \\
CNNBlock(128, dropout 0.3)      & $(16, 16, 128)$ \\
CNNBlock(256, dropout 0.3)      & $(8, 8, 256)$ \\
GlobalAveragePooling2D          & $(256)$ \\
Dense(256, ReLU, L2 $10^{-4}$) + Dropout(0.5) & $(256)$ \\
Dense(8, Softmax)               & $(8)$ \\
\bottomrule
\end{tabular}
\end{center}

Svaki \texttt{CNNBlock} čine: Conv2D $\rightarrow$ BatchNorm $\rightarrow$
ReLU $\rightarrow$ MaxPool2D $\rightarrow$ [Dropout]. Model ima
$\approx$458k parametara.

\subsection{Optimizacija i stabilnost treninga}

Mali i šumovit skup podataka zahteva konzervativnu optimizaciju; sledeća
konfiguracija daje stabilne krive treninga bez divergencije:

\begin{table}[H]
\centering
\begin{tabular}{lll}
\toprule
Parametar & Vrednost & Obrazloženje \\
\midrule
Optimizator       & Adam, $LR = 3\cdot10^{-4}$ & veći LR divergira na ovom problemu \\
Gradient clipping & clipnorm $=1.0$ & seče retke eksplozivne gradijente \\
Batch size        & 64 & mirnije BatchNorm statistike \\
\bottomrule
\end{tabular}
\caption{Konfiguracija optimizacije.}
\end{table}

\subsection{Selekcija modela}

Selekcija najboljeg modela prati validacionu tačnost, a ne validacioni
gubitak: sa focal loss funkcijom na malom validacionom skupu (216 uzoraka)
validacioni gubitak je izrazito šumovit, pa bi selekcija po njemu birala
slučajno povoljne rane epohe. Raspored learning rate-a i dalje prati
validacioni gubitak, jer je za tu svrhu gladak signal dovoljan.

\begin{table}[H]
\centering
\begin{tabular}{ll}
\toprule
Callback & Parametri \\
\midrule
EarlyStopping     & monitor $=$ val\_accuracy, patience $= 20$, restore best \\
ModelCheckpoint   & monitor $=$ val\_accuracy, save best only \\
ReduceLROnPlateau & monitor $=$ val\_loss, factor $= 0.5$, patience $= 5$ \\
\bottomrule
\end{tabular}
\caption{Callback konfiguracija.}
\end{table}

\subsection{Funkcija gubitka: focal loss}

\begin{equation}
  \mathrm{FL}(p_t) = -(1 - p_t)^{\gamma}\,\log(p_t), \qquad \gamma = 2{,}0
\end{equation}

Focal loss smanjuje doprinos lako klasifikovanih primera i fokusira trening
na teške klase (\texttt{sad}, \texttt{fearful}).

\begin{figure}[H]
  \centering
  \includegraphics[width=\textwidth]{training_history.png}
  \caption{Krive treninga CNN-a od nule (tačnost i gubitak).}
  \label{fig:hist-cnn}
\end{figure}

\section{Transfer learning: EfficientNetB0}

\subsection{Arhitektura}

\begin{center}
\begin{tabular}{ll}
\toprule
Sloj & Uloga \\
\midrule
Ulaz $(128,128,3)$      & mel + $\Delta$ + $\Delta\Delta$ \\
ZScoreToPixels          & $[-3,3] \rightarrow [0,255]$ (opseg ImageNet piksela) \\
Resizing$(224,224)$     & nativna rezolucija EfficientNet pretreninga \\
EfficientNetB0          & ImageNet težine, bez top sloja \\
GAP + BatchNorm         & agregacija \\
Dense(256) + Dropout(0.5) + Dense(8, Softmax) & klasifikaciona glava \\
\bottomrule
\end{tabular}
\end{center}

Dve ključne adaptacije ulaza:
\begin{itemize}
  \item \textbf{ZScoreToPixels:} EfficientNetB0 je pretreniran na
        pikselima u opsegu $[0,255]$; bez ovog sloja pretrenirani filteri
        vide vrednosti $\sim[-3,3]$ i ne aktiviraju se;
  \item \textbf{Resizing(224):} na ulazu $128\times128$ poslednja mapa
        EfficientNet-a je samo $4\times4$; na $224\times224$ mreža radi na
        skali na kojoj je pretrenirana.
\end{itemize}

\subsection{Dvofazno treniranje}

\textbf{Faza 1, zamrznut base (do 40 epoha, $LR=10^{-3}$):} trenira se samo
klasifikaciona glava, a rano zaustavljanje odlučuje stvarni kraj faze. Cilj
je da glava nauči da koristi pretrenirane feature mape bez uništavanja
pretreniranih težina naglim gradijentima.

\textbf{Faza 2, fine-tuning (do 60 epoha, $LR=5\cdot10^{-5}$):} odmrzava se
poslednjih 20 slojeva EfficientNet-a. Mali broj odmrznutih slojeva i nizak
learning rate čuvaju pretrenirane reprezentacije i sprečavaju
preprilagođavanje na malom skupu; agresivnije postavke (više slojeva, veći
LR) dovode do trenutnog preprilagođavanja.

Da bi na disku uvek ostao najbolji model iz obe faze, checkpoint u fazi 2
nasleđuje najbolju validacionu tačnost iz faze 1 preko parametra
\texttt{initial\_value\_threshold}; lošija epoha faze 2 tako ne može da
pregazi bolji model iz faze 1.

\begin{figure}[H]
  \centering
  \includegraphics[width=\textwidth]{training_history_transfer.png}
  \caption{Krive treninga transfer modela; isprekidana linija označava
  početak faze 2 (fine-tuning).}
  \label{fig:hist-transfer}
\end{figure}

\section{Ansambl}

Finalni klasifikator je prosek softmax izlaza tri mreže: dva CNN-a od nule
(nezavisna trening run-a) i transfer modela. Različite arhitekture greše na
različitim uzorcima, pa usrednjavanje verovatnoća poništava deo
nekoreliranih grešaka.

Kombinacija članova ansambla birana je na validacionom skupu (73{,}2\%
validacione tačnosti), pa je tek onda jednom izmerena na test skupu, čime
je izbegnuto prilagođavanje test rezultatu.

\section{Rezultati i analiza}

\subsection{Tačnost po modelima}

\begin{table}[H]
\centering
\begin{tabular}{lrr}
\toprule
Model & Validacija & Test \\
\midrule
CNN od nule                  & 68{,}9\% & 59{,}7\% \\
Transfer (EfficientNetB0)    & 66{,}7\% & 60{,}4\% \\
\textbf{Ansambl (2$\times$CNN + transfer)} & \textbf{73{,}2\%} & \textbf{66{,}0\%} \\
\bottomrule
\end{tabular}
\caption{Tačnost svih modela. Ansambl od 66{,}0\% je iznad izmerene
ljudske tačnosti ($\approx$62\%) za prepoznavanje emocija iz audio snimka
govora na RAVDESS bazi.}
\end{table}

\subsection{Per-klasa analiza (ansambl)}

\begin{table}[H]
\centering
\begin{tabular}{lrrr}
\toprule
Klasa & Precision & Recall & F1 \\
\midrule
disgust   & 0.89 & 0.84 & \textbf{0.86} \\
angry     & 0.81 & 0.89 & \textbf{0.85} \\
surprised & 0.78 & 0.90 & \textbf{0.84} \\
calm      & 0.50 & 0.74 & 0.60 \\
happy     & 0.80 & 0.42 & 0.55 \\
neutral   & 0.46 & 0.60 & 0.52 \\
fearful   & 0.56 & 0.47 & 0.51 \\
sad       & 0.47 & 0.37 & 0.41 \\
\bottomrule
\end{tabular}
\caption{Per-klasa metrike ansambla (makro F1 $= 0{,}64$). Tri klase su na
F1 $\geq 0{,}84$; najteža je \texttt{sad}: tiha, spora, akustički slična
klasama \texttt{calm} i \texttt{neutral}.}
\end{table}

\begin{figure}[H]
  \centering
  \includegraphics[width=0.9\textwidth]{confusion_matrix.png}
  \caption{Konfuziona matrica ansambla (procenti po redu).}
  \label{fig:cm}
\end{figure}

\subsection{Analiza preprilagođavanja}

Jaz između trening i validacione tačnosti (trening $\sim$90\%, validacija
67--73\%) postoji u svim modelima i posledica je fundamentalnog
ograničenja: $\sim$135 originalnih snimaka po klasi. Augmentacija umnožava
varijacije, ali ne dodaje nove govornike ni nove interpretacije.
Regularizacija (Dropout, L2, SpecAugment, focal loss) usporava memorisanje
trening skupa, ali ga ne može eliminisati.

Jaz između validacije i testa (73\% naspram 66\%) ima dva uzroka: (1)
validacioni skup učestvuje u svim odlukama (selekcija epohe, raspored
learning rate-a, izbor ansambla), pa je procena na njemu blago
optimistična; (2) oba skupa su mala (216 / 144 uzoraka), pa nekoliko
uzoraka pravi procentni poen razlike.

\subsection{Realan plafon zadatka}

\begin{table}[H]
\centering
\begin{tabular}{ll}
\toprule
Ograničenje & Uticaj \\
\midrule
1.440 snimaka, 8 klasa & $\sim$135 originala po klasi, glavni limit \\
Ljudska tačnost $\approx$62\% & glumljene emocije nisu jednoznačne ni ljudima \\
Akustičko preklapanje & sad$\leftrightarrow$calm$\leftrightarrow$neutral;
                        happy$\leftrightarrow$fearful$\leftrightarrow$surprised \\
Test $=$ 144 uzorka & $\pm$3 p.p.\ šuma u svakom merenju \\
\bottomrule
\end{tabular}
\caption{Ograničenja zadatka.}
\end{table}

Radovi koji prijavljuju 75--85\% koriste wav2vec2/HuBERT embeddinge
(pretrenirane na hiljadama sati govora, što izlazi iz okvira pristupa
,,CNN nad spektrogramima``), spajanje više baza (RAVDESS+TESS+CREMA-D) ili
multimodalnost (audio+video). U okviru zadatka ovog projekta, 66\% je
blizu realnog plafona.

\section{Zaključak prvog dela}

Implementiran je kompletan pipeline za prepoznavanje emocija iz govora, od
sirovih WAV fajlova do evaluacije klasifikatora. Ključni elementi rešenja:

\begin{enumerate}
  \item \textbf{Podaci:} trim tišine, delta i delta-delta kanali (dinamika
        govora), tačno balansirana augmentacija (432 uzorka po klasi),
        reprodukovljiv i paralelizovan preprocesing.
  \item \textbf{Trening:} konzervativna optimizacija ($LR = 3\cdot10^{-4}$,
        gradient clipping, batch 64) i selekcija modela po validacionoj
        tačnosti.
  \item \textbf{Transfer learning:} adaptacija ulaza (ZScoreToPixels,
        Resizing na 224) i pažljiv dvofazni fine-tuning (20 slojeva,
        $LR = 5\cdot10^{-5}$).
  \item \textbf{Ansambl:} prosek softmax izlaza tri mreže, biran na
        validacionom skupu.
\end{enumerate}

\textbf{Konačni rezultat: 66{,}0\% tačnosti na test skupu za 8 klasa}
(makro F1 $= 0{,}64$), iznad izmerene ljudske tačnosti za ovaj zadatak
($\approx$62\%). Dalji značajan napredak zahtevao bi pretrenirane govorne
reprezentacije (wav2vec2/HuBERT) ili veće skupove podataka, što izlazi iz
okvira zadatka.

% ============================================================
% DEO 2: FUZZY LOGIKA
% ------------------------------------------------------------
% Ovde dodati drugi deo izveštaja, npr.:
%
% \part{Fuzzy logika}
% \section{...}
%
% ============================================================

\end{document}
